\section*{ВВЕДЕНИЕ}
\addcontentsline{toc}{section}{ВВЕДЕНИЕ}

В последнее десятилетие в мировой экономике наблюдается стремительный рост модели потребления, основанной на совместном использовании товаров, получившей название «шеринг-экономика» (от англ. sharing — делить). Идея временного пользования вещью вместо её приобретения в собственность становится всё более популярной среди населения. Аренда электроинструментов, бытовой техники, спортивного инвентаря, одежды и даже транспортных средств позволяет потребителям существенно экономить ресурсы, а владельцам вещей — получать прибыль от простаивающего имущества. Данная модель потребления не только выгодна экономически, но и отвечает современным принципам устойчивого развития, снижая нагрузку на экологию за счёт уменьшения объёмов производства и утилизации товаров.

Развитие информационных технологий и повсеместное проникновение сети Интернет стали ключевыми факторами, обеспечившими рост рынка арендных услуг. Если раньше аренда ограничивалась специализированными пунктами проката, расположенными в конкретных районах города, то сегодня веб-сайты и мобильные приложения позволяют создать единое цифровое пространство, где арендодатели и арендаторы могут быстро находить друг друга независимо от времени и местоположения.

Сайт для аренды вещей выполняет функции не только витрины товаров, но и полноценного посреднического сервиса. Он автоматизирует процессы бронирования, расчёта стоимости аренды на срок, обработки платежей и коммуникации между сторонами. Для современного бизнеса в сфере проката наличие качественного веб-сайта является необходимым условием конкурентоспособности. От того, насколько удобен и функционален интерфейс, зависит, сможет ли компания привлечь и удержать клиентов в условиях высокой конкуренции с крупными маркетплейсами и специализированными агрегаторами.

Клиентская часть сайта (фронтенд) играет в этом процессе ключевую роль. Именно через интерфейс пользователь взаимодействует с сервисом: просматривает каталог, изучает характеристики вещей, выбирает даты аренды и совершает заказ. От качества вёрстки, скорости загрузки, адаптивности под мобильные устройства и интуитивной понятности кнопок и форм напрямую зависит, превратится ли посетитель в реального клиента. Современные веб-технологии позволяют создавать интерфейсы, которые работают с серверной частью сайта в фоновом режиме (асинхронно), благодаря чему оформление аренды происходит без задержек и перезагрузки страниц, максимально приближаясь по удобству к настольным приложениям.

\emph{Цель настоящей работы} – разработка клиентской части веб-сайта для сервиса аренды вещей, включая создание пользовательского интерфейса и реализацию программного взаимодействия с серверным API. Для достижения поставленной цели необходимо решить \emph{следующие задачи:}
\begin{itemize}
\item провести анализ предметной области сервисов аренды и существующих решений;
\item разработать структуру и дизайн пользовательского интерфейса;
\item выполнить вёрстку страниц сайта с использованием современных технологий HTML и CSS;
\item реализовать на языке JavaScript интерактивные элементы и функции обращения к API для выполнения операций (получение списка вещей, оформление заказа, фильтрация);
\item провести тестирование разработанного интерфейса и его интеграцию с серверной частью.
\end{itemize}

\emph{Структура и объем работы.} Отчет состоит из введения, 4 разделов основной части, заключения, списка использованных источников, 2 приложений. Текст выпускной квалификационной работы равен \formbytotal{lastpage}{страниц}{е}{ам}{ам}.

\emph{Во введении} сформулирована цель работы, поставлены задачи разработки, описана структура работы, приведено краткое содержание каждого из разделов.

\emph{В первом разделе} рассматриваются современное состояние рынка арендных онлайн-сервисов, основные требования к подобным сайтам и существующие аналоги. Проводится анализ предметной области и формируется техническое задание на разработку клиентской части.

\emph{Во втором разделе} описываются архитектура разрабатываемого интерфейса, структура страниц, проектируются основные пользовательские сценарии (поиск вещи, оформление аренды). Разрабатываются прототипы интерфейса и схемы взаимодействия с API.

\emph{В третьем разделе} подробно описывается процесс реализации клиентской части: создание HTML-разметки, написание стилей CSS, разработка логики на JavaScript, реализация вызовов к API для отправки и получения данных. Приводятся примеры кода.

\emph{В четвертом разделе} описываются результаты тестирования разработанного интерфейса в различных браузерах и на различных устройствах, проверяется корректность работы с API, оценивается производительность.

В заключении излагаются основные результаты работы, полученные в ходе разработки, и делаются выводы о соответствии готового продукта поставленным задачам.

В приложении А представлены экранные формы (скриншоты) разработанного сайта. 
В приложении Б представлены фрагменты исходного кода (HTML, CSS, JavaScript).
